\documentclass[../main.tex]{subfiles}
\begin{document}

\section*{Bối cảnh}
\addcontentsline{toc}{section}{Bối cảnh}

Trong những năm gần đây, bức tranh an ninh mạng trên hệ điều hành Windows đã có sự thay đổi đáng kể cả về quy mô lẫn mức độ tinh vi của các cuộc tấn công. Thay vì sử dụng các mã độc truyền thống dễ bị nhận diện, các tác nhân đe dọa ngày càng ưu tiên những chiến thuật khó phát hiện hơn như Advanced Persistent Threat (APT), trong đó kẻ tấn công duy trì sự hiện diện lâu dài và thực hiện các bước xâm nhập một cách âm thầm. Đáng chú ý, nhiều chiến dịch hiện đại tận dụng kỹ thuật Living-off-the-Land, khai thác các công cụ hợp pháp sẵn có trong hệ thống để thực thi hành vi độc hại, qua đó làm suy giảm đáng kể hiệu quả của các cơ chế phát hiện dựa trên chữ ký.

Song song với sự gia tăng về độ tinh vi của tấn công là sự bùng nổ về khối lượng dữ liệu giám sát. Các hệ thống như Sysmon có thể ghi nhận hàng trăm nghìn đến hàng chục triệu sự kiện mỗi ngày trên một máy, tạo ra áp lực rất lớn đối với các trung tâm điều hành an ninh (SOC). Trong thực tế, lượng dữ liệu này không chỉ gây khó khăn cho việc lưu trữ mà còn làm giảm hiệu quả của quá trình phân tích, khi các chuyên gia phải xử lý một số lượng lớn cảnh báo nhưng chỉ một phần nhỏ trong đó thực sự liên quan đến hành vi tấn công.

Báo cáo năm 2023 của IBM Security cho thấy thời gian trung bình để phát hiện và xử lý một sự cố an ninh lên tới 277 ngày. Con số này phản ánh một khoảng cách đáng kể giữa khả năng phát hiện tự động và hiệu quả điều tra thực tế, đặc biệt trong giai đoạn triage ban đầu. Trong bối cảnh đó, việc phát triển các phương pháp phân tích log hiệu quả, có khả năng giảm tải cho SOC đồng thời nâng cao chất lượng phát hiện, trở thành một yêu cầu cấp thiết.

\section*{Hạn chế của các giải pháp triage hiện tại}
\addcontentsline{toc}{section}{Hạn chế của các giải pháp triage hiện tại}

Mặc dù đã có nhiều công cụ hỗ trợ phân tích log Windows, các giải pháp hiện tại vẫn tồn tại những hạn chế mang tính hệ thống. Một trong những vấn đề cốt lõi là sự phụ thuộc mạnh vào các luật hoặc chữ ký được định nghĩa sẵn. Cách tiếp cận này vốn hiệu quả đối với các kỹ thuật đã được biết đến, nhưng lại tỏ ra kém thích ứng trước các biến thể mới hoặc các kịch bản tấn công sử dụng LOLBins, nơi hành vi độc hại được che giấu dưới hình thức hợp lệ.

Bên cạnh đó, nhiều hệ thống áp dụng các kỹ thuật pruning (cắt tỉa) đồ thị tiến trình nhằm giảm nhiễu dữ liệu và cải thiện hiệu năng xử lý. Tuy nhiên, việc loại bỏ các node có tần suất thấp thường dẫn đến hệ quả không mong muốn là làm mất đi các bằng chứng pháp y quan trọng. Các lệnh trinh sát như \texttt{whoami} hay \texttt{ipconfig}, vốn chỉ xuất hiện một lần trong chuỗi tấn công, có thể bị loại bỏ, khiến cho cấu trúc phả hệ tiến trình bị đứt gãy và làm suy giảm khả năng tái dựng toàn bộ kịch bản xâm nhập.

Một hạn chế khác nằm ở yêu cầu về dữ liệu nền (baseline). Nhiều phương pháp phát hiện bất thường cần một tập dữ liệu "sạch" để huấn luyện hoặc so sánh, trong khi điều này rất khó đảm bảo trong môi trường thực tế, nơi hành vi bình thường và bất thường thường xen kẽ nhau. Ngoài ra, các phương pháp học máy phổ biến như Isolation Forest hay OneClass SVM thường hoạt động theo cơ chế khó giải thích (black-box). Cuối cùng, việc sử dụng các ngưỡng cứng hoặc cơ chế cộng điểm tuyến tính trong nhiều hệ thống dẫn đến hiện tượng "vùng chết toán học" (Dead Zone), trong đó các tín hiệu bất thường mức trung bình không đủ mạnh để vượt qua ngưỡng và do đó bị bỏ qua hoàn toàn.

\section*{Khoảng trống cần lấp}
\addcontentsline{toc}{section}{Khoảng trống cần lấp}

Từ những hạn chế trên, có thể nhận thấy sự thiếu hụt của một phương pháp triage vừa linh hoạt, vừa đảm bảo tính toàn vẹn của dữ liệu pháp y. Trước hết, cần có một hệ thống có khả năng hoạt động độc lập (offline), không phụ thuộc vào dữ liệu huấn luyện bên ngoài, nhưng vẫn có thể tự xây dựng baseline từ chính dữ liệu đang phân tích (self-learning). Đồng thời, việc bảo toàn đầy đủ cấu trúc đồ thị tiến trình, bao gồm cả các nhánh chỉ xuất hiện một lần (unique branches), là yêu cầu thiết yếu để đảm bảo khả năng tái dựng chuỗi tấn công (attack chain). Quan trọng hơn, hệ thống cần cung cấp một cơ chế định lượng rủi ro có thể giải thích được (Explainable AI — XAI), cùng một cách tiếp cận toán học phù hợp để tránh hiện tượng triệt tiêu tín hiệu (Dead Zone).

\section*{Mục tiêu đề tài}
\addcontentsline{toc}{section}{Mục tiêu đề tài}

Trên cơ sở đó, đề tài này hướng tới việc xây dựng một hệ thống triage log Windows theo hướng không phụ thuộc vào danh sách trắng hay chữ ký, đồng thời đảm bảo khả năng giải thích của toàn bộ quá trình phân tích. Mục tiêu cụ thể là phát triển một công cụ có thể phát hiện các hành vi bất thường trong log Sysmon với độ bao phủ cao, hướng tới đạt mức recall trên 90\% khi đánh giá trên các tập dữ liệu tấn công công khai.

\section*{Phạm vi nghiên cứu}
\addcontentsline{toc}{section}{Phạm vi nghiên cứu}

Phạm vi của đề tài được giới hạn ở việc phân tích log sự kiện Windows thu thập từ Sysmon, với định dạng \texttt{.evtx} và tập trung vào một số nhóm sự kiện quan trọng liên quan đến tiến trình, mạng, và thay đổi hệ thống. Hệ thống được thiết kế để hoạt động trong môi trường Windows hiện đại, bao gồm Windows 10, Windows 11 và các phiên bản Windows Server tương đương, chỉ xem xét kịch bản phân tích offline trên một máy đơn lẻ.

\section*{Phương pháp nghiên cứu}
\addcontentsline{toc}{section}{Phương pháp nghiên cứu}

Đề tài áp dụng phương pháp nghiên cứu thực nghiệm nhằm đánh giá hiệu quả của hệ thống sau khi triển khai. Cụ thể, kết quả của hệ thống sẽ được so sánh với các công cụ triage phổ biến như Hayabusa và Chainsaw trên cùng một tập dữ liệu. Các chỉ số định lượng như recall, precision, F1-score, hiệu năng xử lý và mức sử dụng tài nguyên sẽ được sử dụng để đánh giá.

\section*{Đóng góp dự kiến}
\addcontentsline{toc}{section}{Đóng góp dự kiến}

Đề tài kỳ vọng đóng góp một cách tiếp cận mới trong bài toán triage log Windows, trong đó việc giảm nhiễu dữ liệu được thực hiện mà không làm mất đi các bằng chứng quan trọng. Đồng thời, hệ thống đề xuất một cơ chế tự học baseline từ dữ liệu đầu vào (on-the-fly self-learning), kết hợp với các phương pháp thống kê và mô hình hóa hành vi nhằm phát hiện bất thường một cách toàn diện. Ngoài ra, việc tích hợp cơ chế định lượng rủi ro có thể giải thích (explainable risk scoring), cùng với phương pháp lan truyền rủi ro trong đồ thị tiến trình (risk propagation), giúp nâng cao khả năng phát hiện và hỗ trợ hiệu quả cho quá trình điều tra.

\section*{Định vị hệ thống SRAH như một môi trường thử nghiệm (Testbed)}
\addcontentsline{toc}{section}{Định vị hệ thống SRAH như một môi trường thử nghiệm (Testbed)}

Một điểm cần đặc biệt lưu ý là hệ thống SRAH (Statistical Anomaly \& Dynamic Guilt Backpropagation) được đề xuất trong đồ án này không hướng tới việc xây dựng một sản phẩm phần mềm thương mại hoàn chỉnh hay một hệ thống phòng thủ hoạt động độc lập trong môi trường sản xuất (production system). Thay vào đó, SRAH được thiết kế và định vị như một môi trường thử nghiệm khoa học (scientific testbed) và là một công cụ thực chứng (instrumentation tool).

Mục đích chính của SRAH là cung cấp một nền tảng thực nghiệm để triển khai và đo lường một cách định lượng giới hạn của các phương pháp phát hiện bất thường dựa trên thống kê không giám sát (unsupervised statistical anomaly detection) trên nhật ký sự kiện Windows. Việc định vị này cho phép chúng ta cô lập các biến số hành vi, thực hiện các nghiên cứu loại trừ (ablation studies) để đánh giá độ nhạy và hiệu quả đóng góp của từng thành phần toán học (Point, Contextual, Collective, Sequential Anomalies) đối với từng lớp tấn công cụ thể. Do đó, các kết quả thực nghiệm đạt được sẽ phản ánh một cách khách quan tiềm năng và giới hạn thực tế của các lý thuyết toán học trong lĩnh vực điều tra số và ứng phó sự cố (DFIR), thay vì chỉ đánh giá khả năng tối ưu hóa kỹ thuật của một sản phẩm thương mại.

\section*{Bố cục đồ án}
\addcontentsline{toc}{section}{Bố cục đồ án}

Phần còn lại của đồ án được tổ chức như sau:

\begin{itemize}
    \item \textbf{Chương 1: Kiến thức cơ sở} Trình bày các kiến thức nền tảng và cơ sở lý thuyết liên quan đến phát hiện bất thường (anomaly detection) và phân tích nhật ký (log analysis), bao gồm toàn bộ cơ sở toán học của các thành phần scoring, các công trình nghiên cứu liên quan và phân tích khoảng trống nghiên cứu.
    \item \textbf{Chương 2: Thiết kế hệ thống} Mô tả chi tiết kiến trúc và các tầng xử lý của hệ thống SRAH từ Layer 0 đến Layer 4, bao gồm việc lựa chọn công nghệ PyO3 và thiết lập các tham số cụ thể dựa trên cơ sở lý thuyết và thực nghiệm.
    \item \textbf{Chương 3: Thực nghiệm và đánh giá} Tập trung vào thiết kế thực nghiệm, mô tả dữ liệu kiểm thử, trình bày và phân tích chi tiết kết quả thử nghiệm định lượng của hệ thống SRAH cùng các nghiên cứu loại trừ (ablation study), đánh giá độ nhạy tham số và thảo luận.
    \item \textbf{Kết luận và hướng phát triển} Tổng kết các kết quả đạt được của đồ án, chỉ ra những hạn chế hiện tại và đề xuất các hướng nghiên cứu tiếp theo nhằm hoàn thiện giải pháp.
\end{itemize}

\end{document}